%%
%% This is file `sample-sigconf.tex',
%% generated with the docstrip utility.
%%
%% The original source files were:
%%
%% samples.dtx  (with options: `all,proceedings,bibtex,sigconf')
%% 
%% IMPORTANT NOTICE:
%% 
%% For the copyright see the source file.
%% 
%% Any modified versions of this file must be renamed
%% with new filenames distinct from sample-sigconf.tex.
%% 
%% For distribution of the original source see the terms
%% for copying and modification in the file samples.dtx.
%% 
%% This generated file may be distributed as long as the
%% original source files, as listed above, are part of the
%% same distribution. (The sources need not necessarily be
%% in the same archive or directory.)
%%
%%
%% Commands for TeXCount
%TC:macro \cite [option:text,text]
%TC:macro \citep [option:text,text]
%TC:macro \citet [option:text,text]
%TC:envir table 0 1
%TC:envir table* 0 1
%TC:envir tabular [ignore] word
%TC:envir displaymath 0 word
%TC:envir math 0 word
%TC:envir comment 0 0
%%
%% The first command in your LaTeX source must be the \documentclass
%% command.
%%
%% For submission and review of your manuscript please change the
%% command to \documentclass[manuscript, screen, review]{acmart}.
%%
%% When submitting camera ready or to TAPS, please change the command
%% to \documentclass[sigconf]{acmart} or whichever template is required
%% for your publication.
%%
%%
\documentclass[sigconf]{acmart}
%%
%% \BibTeX command to typeset BibTeX logo in the docs
\AtBeginDocument{%
  \providecommand\BibTeX{{%
    Bib\TeX}}}

%% Rights management information.  This information is sent to you
%% when you complete the rights form.  These commands have SAMPLE
%% values in them; it is your responsibility as an author to replace
%% the commands and values with those provided to you when you
%% complete the rights form.
\setcopyright{acmlicensed}
\copyrightyear{2018}
\acmYear{2018}
\acmDOI{XXXXXXX.XXXXXXX}
%% These commands are for a PROCEEDINGS abstract or paper.
\acmConference[Conference acronym 'XX]{Make sure to enter the correct
  conference title from your rights confirmation email}{June 03--05,
  2018}{Woodstock, NY}
%%
%%  Uncomment \acmBooktitle if the title of the proceedings is different
%%  from ``Proceedings of ...''!
%%
%%\acmBooktitle{Woodstock '18: ACM Symposium on Neural Gaze Detection,
%%  June 03--05, 2018, Woodstock, NY}
\acmISBN{978-1-4503-XXXX-X/2018/06}


%%
%% Submission ID.
%% Use this when submitting an article to a sponsored event. You'll
%% receive a unique submission ID from the organizers
%% of the event, and this ID should be used as the parameter to this command.
%%\acmSubmissionID{123-A56-BU3}

%%
%% For managing citations, it is recommended to use bibliography
%% files in BibTeX format.
%%
%% You can then either use BibTeX with the ACM-Reference-Format style,
%% or BibLaTeX with the acmnumeric or acmauthoryear sytles, that include
%% support for advanced citation of software artefact from the
%% biblatex-software package, also separately available on CTAN.
%%
%% Look at the sample-*-biblatex.tex files for templates showcasing
%% the biblatex styles.
%%

%%
%% The majority of ACM publications use numbered citations and
%% references.  The command \citestyle{authoryear} switches to the
%% "author year" style.
%%
%% If you are preparing content for an event
%% sponsored by ACM SIGGRAPH, you must use the "author year" style of
%% citations and references.
%% Uncommenting
%% the next command will enable that style.
%%\citestyle{acmauthoryear}


% \usepackage{amsmath, mathtools, amssymb, amsthm}
\usepackage{amsmath, mathtools, amsthm}
\usepackage{proof}
\usepackage{xcolor}
\usepackage{subfiles}
\usepackage{stmaryrd}
\usepackage{hyperref}
\usepackage{tcolorbox}
\usepackage{tabularx}
\usepackage{pifont}
\usepackage{enumitem}
\usepackage{threeparttable}
\usepackage{multicol}

\tcbuselibrary{skins}

\definecolor{xhs_blue}{HTML}{364FC7}
\definecolor{xhs_blue_dark}{HTML}{081C59}
\definecolor{xhs_blue_light}{HTML}{EDF2FF}

\definecolor{xhs_red}{HTML}{c13c33}
\definecolor{xhs_green}{HTML}{3c5840}

\definecolor{xhs_all}{HTML}{9f1f31}
\definecolor{xhs_ex}{HTML}{2885a9}


\NewTotalTColorBox{\BlueQuote}{ +m +m }{ 
    notitle,
    colback=xhs_blue_light,
    colbacklower=white,
    frame hidden,
    boxrule=0pt,
    bicolor,
    sharp corners,
    borderline west={2.5pt}{0pt}{xhs_blue}
}{
    \textcolor{xhs_blue}{ \textbf{#1} }
    #2 
}
\renewcommand{\subsectionautorefname}{Section} 

\newcommand{\funcText}[1]{{\ttfamily #1}}

\newcommand{\ruleNameText}[1]{{\ttfamily #1}}

\newcommand{\TFact}[1]{{\ttfamily #1}}
\newcommand{\TLemma}[1]{{\ttfamily #1}}
\newcommand{\TEmpFact}[1]{\lbrack \  #1 \rbrack}
\newcommand{\TFacts}[1]{\lbrack\ #1 \ \rbrack}
\newcommand{\TEmpActFact}[0]{-\hspace{-0.45em}\rightarrow}
\newcommand{\TEmpActFactSL}[1]{- \hspace{-2.3pt} \TEmpFact{#1} \hspace{-4.8pt} \rightarrow} % single line
\newcommand{\TActFacts}[1]{- \hspace{-4.5pt} \TFacts{#1} \hspace{-4.8pt} \rightarrow}
\newcommand{\TActFactsSL}[1]{- \hspace{-2.3pt} \TFacts{#1} \hspace{-4.8pt} \rightarrow}
\newcommand{\TPair}[1]{\langle #1 \rangle}
\newcommand{\TPairN}[1]{$\langle$\text{\ttfamily #1}$\rangle$}

\newcommand{\TLet}[0]{\text{\Tterm{let}}}
\newcommand{\TIn}[0]{\text{\Tterm{in}}}

\newcommand{\TLetS}[0]{\text{\Tterm{let}} \hspace{0.5em}}
\newcommand{\TInS}[0]{\text{\Tterm{in}} \hspace{0.5em}}

\newcommand{\TAll}[0]{\text{\textcolor{xhs_all}{\TLemma{All}}}}
\newcommand{\TEx}[0]{\text{\textcolor{xhs_ex}{\TLemma{Ex}}}}
\newcommand{\TNot}[0]{\text{\TLemma{not}}}

\newcommand{\Tterm}[1]{{\ttfamily #1}}
\newcommand{\Tequation}[1]{{\ttfamily #1}}
\newcommand{\Tfunction}[1]{{\ttfamily #1}}
\newcommand{\Tamarin}[0]{\textsc{Tamarin}}

\newcommand{\TFrTerm}[1]{{\Tterm{$\sim$#1}}}
\newcommand{\TFactTerm}[2]{{\ttfamily #1}$(\,$\Tterm{#2}$\,)$}
\newcommand{\TFrFactTerm}[1]{{\ttfamily Fr}$(\,$\TFrTerm{#1}$\,)$}

\newcommand{\fieldtext}[1]{\texttt{#1}}
\newcommand{\csttext}[1]{\textbf{#1}}
\newcommand{\functext}[1]{\textsf{#1}}

\hbadness=10000

\newcolumntype{S}{>{\centering\arraybackslash}m{3cm}@{*}}

%%
%% end of the preamble, start of the body of the document source.
\begin{document}

%%
%% The "title" command has an optional parameter,
%% allowing the author to define a "short title" to be used in page headers.
\title{ Small Leaks Sink Great Ships: \\Automated Analysis of Protocols that Tolerate Low-Entropy Keys }

% %%
% %% The "author" command and its associated commands are used to define
% %% the authors and their affiliations.
% %% Of note is the shared affiliation of the first two authors, and the
% %% "authornote" and "authornotemark" commands
% %% used to denote shared contribution to the research.
% \author{Ben Trovato}
% \authornote{Both authors contributed equally to this research.}
% \email{trovato@corporation.com}
% \orcid{1234-5678-9012}
% \author{G.K.M. Tobin}
% \authornotemark[1]
% \email{webmaster@marysville-ohio.com}
% \affiliation{%
%   \institution{Institute for Clarity in Documentation}
%   \city{Dublin}
%   \state{Ohio}
%   \country{USA}
% }

% \author{Lars Th{\o}rv{\"a}ld}
% \affiliation{%
%   \institution{The Th{\o}rv{\"a}ld Group}
%   \city{Hekla}
%   \country{Iceland}}
% \email{larst@affiliation.org}

% \author{Valerie B\'eranger}
% \affiliation{%
%   \institution{Inria Paris-Rocquencourt}
%   \city{Rocquencourt}
%   \country{France}
% }

% \author{Aparna Patel}
% \affiliation{%
%  \institution{Rajiv Gandhi University}
%  \city{Doimukh}
%  \state{Arunachal Pradesh}
%  \country{India}}

% \author{Huifen Chan}
% \affiliation{%
%   \institution{Tsinghua University}
%   \city{Haidian Qu}
%   \state{Beijing Shi}
%   \country{China}}

% \author{Charles Palmer}
% \affiliation{%
%   \institution{Palmer Research Laboratories}
%   \city{San Antonio}
%   \state{Texas}
%   \country{USA}}
% \email{cpalmer@prl.com}

% \author{John Smith}
% \affiliation{%
%   \institution{The Th{\o}rv{\"a}ld Group}
%   \city{Hekla}
%   \country{Iceland}}
% \email{jsmith@affiliation.org}

% \author{Julius P. Kumquat}
% \affiliation{%
%   \institution{The Kumquat Consortium}
%   \city{New York}
%   \country{USA}}
% \email{jpkumquat@consortium.net}

% %%
% %% By default, the full list of authors will be used in the page
% %% headers. Often, this list is too long, and will overlap
% %% other information printed in the page headers. This command allows
% %% the author to define a more concise list
% %% of authors' names for this purpose.
% \renewcommand{\shortauthors}{Trovato et al.}

%%
%% The abstract is a short summary of the work to be presented in the
%% article.
\begin{abstract}
Entropy is crucial for cryptographic security as it prevents adversaries from brute-forcing keys.
However, modern protocols like WPA2-PSK and Bluetooth usually tolerate low-entropy keys for user-friendliness or compatibility.
% ! Rewrite following
The protocol's tolerance for using keys susceptible to brute force attacks by adversaries introduces significant potential security risks. 
Indeed, several real-world attacks have successfully exploited these vulnerabilities.
Current symbolic model analysis tools do not emphasize the strength of keys within protocols, resulting in a limitation where these tools are unable to generally identify attacks associated with low-entropy keys.
This paper introduced a general and automated methodology for modeling low-entropy keys and capturing brute-force attacks.
We developed EntropyVerif, an extension of the symbolic analysis tool \Tamarin{}, which allows users to specify the entropy of keys and automatically detect brute-force attacks.
% Comprehensive case studies were conducted using our EntropyVerif tool to analyze real-world protocols that are tolerant to low-entropy keys. 
Through comprehensive case studies using EntropyVerif, we not only rediscovered known attacks but also uncovered a novel attack on Bluetooth Classic devices, which is termed the Unilateral Multi-Downgrade (UMD) attack. 
This newly identified attack significantly reduces the complexity for adversaries attempting to compromise session keys, thereby enhancing the efficiency of such attacks by a factor of $2^{116}$. 
The results from our case studies highlight the effectiveness of our approach in identifying potential vulnerabilities in protocols that tolerate low-entropy keys.

\end{abstract}

%%
%% The code below is generated by the tool at http://dl.acm.org/ccs.cfm.
%% Please copy and paste the code instead of the example below.
%%
\begin{CCSXML}
<ccs2012>
 <concept>
  <concept_id>00000000.0000000.0000000</concept_id>
  <concept_desc>Do Not Use This Code, Generate the Correct Terms for Your Paper</concept_desc>
  <concept_significance>500</concept_significance>
 </concept>
 <concept>
  <concept_id>00000000.00000000.00000000</concept_id>
  <concept_desc>Do Not Use This Code, Generate the Correct Terms for Your Paper</concept_desc>
  <concept_significance>300</concept_significance>
 </concept>
 <concept>
  <concept_id>00000000.00000000.00000000</concept_id>
  <concept_desc>Do Not Use This Code, Generate the Correct Terms for Your Paper</concept_desc>
  <concept_significance>100</concept_significance>
 </concept>
 <concept>
  <concept_id>00000000.00000000.00000000</concept_id>
  <concept_desc>Do Not Use This Code, Generate the Correct Terms for Your Paper</concept_desc>
  <concept_significance>100</concept_significance>
 </concept>
</ccs2012>
\end{CCSXML}

\ccsdesc[500]{Do Not Use This Code~Generate the Correct Terms for Your Paper}
\ccsdesc[300]{Do Not Use This Code~Generate the Correct Terms for Your Paper}
\ccsdesc{Do Not Use This Code~Generate the Correct Terms for Your Paper}
\ccsdesc[100]{Do Not Use This Code~Generate the Correct Terms for Your Paper}

%%
%% Keywords. The author(s) should pick words that accurately describe
%% the work being presented. Separate the keywords with commas.
\keywords{Do, Not, Use, This, Code, Put, the, Correct, Terms, for,
  Your, Paper}
%% A "teaser" image appears between the author and affiliation
%% information and the body of the document, and typically spans the
%% page.
% \begin{teaserfigure}
%   \includegraphics[width=\textwidth]{sampleteaser}
%   \caption{Seattle Mariners at Spring Training, 2010.}
%   \Description{Enjoying the baseball game from the third-base
%   seats. Ichiro Suzuki preparing to bat.}
%   \label{fig:teaser}
% \end{teaserfigure}

\received{20 February 2007}
\received[revised]{12 March 2009}
\received[accepted]{5 June 2009}

%%
%% This command processes the author and affiliation and title
%% information and builds the first part of the formatted document.
\maketitle


\subfile{sections/section1.tex}
\subfile{sections/section2.tex}
\subfile{sections/section3.tex}
\subfile{sections/section4.tex}
\subfile{sections/section5.tex}
\subfile{sections/section6.tex}
\subfile{sections/section7.tex}
\subfile{sections/section8.tex}

\bibliographystyle{ACM-Reference-Format}
\bibliography{library}

\end{document}
\endinput
%%
%% End of file `sample-sigconf.tex'.
